\SourceLeaf{86}%
\Article{13}
Vergleicht man die Beispiele des vorigen Artikels miteinander, so scheint sich zumindest auf den ersten Blick kein einfaches Kriterium anzubieten, das die ersten Moduln von den zweiten unterscheidet. Dennoch gibt es \emph{zwei} solche Kriterien, die sich durch besondere Eleganz und Einfachheit auszeichnen. Die folgenden Überlegungen werden den Weg zu einem von ihnen bahnen.

Der Modul $p$ lässt sich als Primzahl der Form $8n+1$ auf genau eine Weise in der Form $aa+2bb$ darstellen (\emph{Disquiss. Arithm.}, Artikel 182, II). Dabei setzen wir voraus, dass die Wurzeln $a,b$ positiv genommen werden. Offenbar ist $a$ ungerade und $b$ gerade; wir setzen $b=2^\lambda c$, wobei $c$ ungerade ist. Nun bemerken wir:

I. Da $p\equiv aa\ (\mathrm{mod.}\ c)$ gilt, ist $p$ ein quadratischer Rest von $c$ und somit auch von jedem der Primfaktoren, in die $c$ zerlegt wird. Umgekehrt sind nach dem Fundamentaltheorem diese einzelnen Primfaktoren quadratische Reste von $p$, und damit ist auch ihr Produkt $c$ ein quadratischer Rest von $p$. Da dies auch für die Zahl $2$ gilt, ist $b$ offenbar ein quadratischer Rest von $p$, und somit sind $bb$ und auch $-bb$ biquadratische Reste.

II. Daher muss $-2bb$ zu derselben Klasse gehören wie die Zahl $2$. Da $aa\equiv-2bb$ gilt, liegt $2$ offenbar in der Klasse $A$ oder in der Klasse $C$, je nachdem $a$ ein quadratischer Rest von $p$ oder ein quadratischer Nichtrest ist.

III. Nehmen wir nun an, $a$ sei in seine Primfaktoren zerlegt. Diejenigen der Form $8m+1$ oder $8m+7$ seien mit $\alpha,\alpha',\alpha''$ usw. bezeichnet, diejenigen der Form $8m+3$ oder $8m+5$ mit $\beta,\beta',\beta''$ usw.; die Anzahl der letzteren sei $=\mu$. Da $p\equiv2bb\ (\mathrm{mod.}\ a)$ gilt, ist $p$ ein quadratischer Rest derjenigen Primfaktoren von $a$, deren quadratischer Rest $2$ ist, also der Faktoren $\alpha,\alpha',\alpha''$ usw.; dagegen ist es ein quadratischer Nichtrest derjenigen Faktoren, deren quadratischer Nichtrest $2$ ist, also der Faktoren $\beta,\beta',\beta''$ usw. Umgekehrt sind daher nach dem Fundamentaltheorem die einzelnen $\alpha,\alpha',\alpha''$ usw. quadratische Reste von $p$, die einzelnen $\beta,\beta',\beta''$ usw. hingegen quadratische Nichtreste. Wir schließen also, dass das Produkt $a$ ein quadratischer Rest von $p$ oder ein Nichtrest ist, je nachdem $\mu$ gerade oder ungerade ist.

IV. Man bestätigt leicht, dass das Produkt aller $\alpha,\alpha',\alpha''$ usw. die Form $8m+1$ oder $8m+7$ hat und dass dies auch für das Produkt aller $\beta,\beta',\beta''$ usw. gilt, wenn deren Anzahl gerade ist. In diesem Fall muss also auch das Produkt $a$ die Form $8m+1$ oder $8m+7$ haben. Ist dagegen beim Produkt aller $\beta,\beta',\beta''$ usw.\space%
\SourceLeaf{87}%
die Anzahl der Faktoren ungerade, so hat es die Form $8m+3$ oder $8m+5$, und dasselbe gilt in diesem Fall auch für das Produkt $a$.

Aus alledem ergibt sich der elegante Satz:

\emph{Hat $a$ die Form $8m+1$ oder $8m+7$, so liegt die Zahl $2$ in $A$; hat dagegen $a$ die Form $8m+3$ oder $8m+5$, so liegt die Zahl $2$ in $C$.}

\Needspace{11\baselineskip}\begingroup\rightskip=0pt plus1fil\relax Dies bestätigen die Beispiele des vorigen Artikels. Die ersten Moduln werden nämlich folgendermaßen zerlegt: \mbox{$73=1+2\pdot36$}, \mbox{$89=81+2\pdot4$}, \mbox{$113=81+2\pdot16$}, \mbox{$233=225+2\pdot4$}, \mbox{$257=225+2\pdot16$}, \mbox{$281=81+2\pdot100$}, \mbox{$337=49+2\pdot144$}, \mbox{$353=225+2\pdot64$}; die zweiten dagegen folgendermaßen: \mbox{$17=9+2\pdot4$}, \mbox{$41=9+2\pdot16$}, \mbox{$97=25+2\pdot36$}, \mbox{$137=9+2\pdot64$}, \mbox{$193=121+2\pdot36$}, \mbox{$241=169+2\pdot36$}, \mbox{$313=25+2\pdot144$}, \mbox{$401=9+2\pdot196$}, \mbox{$409=121+2\pdot144$}, \mbox{$433=361+2\pdot36$}, \mbox{$449=441+2\pdot4$}, \mbox{$457=169+2\pdot144$}.\par\endgroup

\Article{14}
Die Zerlegung der Zahl $p$ in ein Quadrat und das Doppelte eines Quadrats hat einen so bemerkenswerten Zusammenhang mit der Einteilung der Zahl $2$ ergeben, dass es lohnend scheint zu untersuchen, ob die Zerlegung in zwei Quadrate, die bekanntlich auch für $p$ möglich ist, vielleicht zu einem ähnlichen Ergebnis führt. Hier stehen daher die Zerlegungen der Zahlen $p$, für die $2$ zur folgenden Klasse gehört:
\[
\begin{array}{c|c}
A&C\\\hline
9+64&1+16\\
25+64&25+16\\
49+64&81+16\\
169+64&121+16\\
1+256&49+144\\
25+256&225+16\\
81+256&169+144\\
289+64&1+400\\
&9+400\\
&289+144\\
&49+400\\
&441+16
\end{array}
\]

\SourceLeaf{88}%
Zunächst bemerken wir, dass von den beiden Quadraten, in die $p$ zerlegt wird, eines ungerade sein muss, das wir $=aa$ setzen, und das andere gerade, das wir $=bb$ setzen. Da $aa$ die Form $8n+1$ hat, entsprechen den Werten von $b$, die das Doppelte einer ungeraden Zahl sind, offenbar Werte von $p$ der Form $8n+5$. Diese bleiben hier bei unserer induktiven Untersuchung ausgeschlossen, da für sie die Zahl $2$ in der Klasse $B$ oder $D$ läge. Für Werte von $p$ der Form $8n+1$ muss hingegen $b$ durch vier teilbar sein. Wenn wir der induktiven Beobachtung vertrauen dürfen, die die angegebene Tabelle vor Augen führt, so gehört die Zahl $2$ zur Klasse $A$ für alle Moduln, bei denen $b$ die Form $8n$ hat, und zur Klasse $C$ für alle Moduln, bei denen $b$ die Form $8n+4$ hat. Dieser Satz verlangt jedoch eine sehr viel tiefere Untersuchung als derjenige, den wir im vorigen Artikel gewonnen haben. Seinem Beweis müssen mehrere vorbereitende Untersuchungen darüber vorausgehen, in welcher Reihenfolge die Zahlen der Sammlungen $A,B,C,D$ aufeinanderfolgen.

\Article{15}
Wir bezeichnen die Anzahl der Zahlen aus $A$, auf die unmittelbar eine Zahl aus $A,B,C,D$ folgt, entsprechend mit $(00),(01),(02),(03)$; ebenso die Anzahl der Zahlen aus $B$, auf die eine Zahl aus $A,B,C,D$ folgt, entsprechend mit $(10),(11),(12),(13)$. Entsprechend seien in $C$ jeweils $(20),(21),(22),(23)$ Zahlen und in $D$ jeweils $(30),(31),(32),(33)$ Zahlen enthalten, auf die eine Zahl aus $A,B,C,D$ folgt. Wir wollen diese sechzehn Anzahlen a priori bestimmen. Damit die Leser die allgemeinen Überlegungen leichter mit Beispielen vergleichen können, erschien es zweckmäßig, hier die Zahlenwerte der Einträge des Schemas $(S)$
\[
\begin{array}{llll}
(00),&(01),&(02),&(03)\\
(10),&(11),&(12),&(13)\\
(20).&(21),&(22),&(23)\\
(30),&(31),&(32),&(33)
\end{array}
\]
für jeden Modul anzugeben, dessen Klasseneinteilung wir in Artikel 11 angegeben haben.

\[
\begin{array}{c|c|c|c}
\begin{gathered}
p=5\\[3pt]
\begin{array}{l}
0,1,0\mathpunct{.}\,0\\
0,0,0,1\\
0,0,0,0\\
0,0,1,0
\end{array}
\end{gathered}&
\begin{gathered}
p=13\\[3pt]
\begin{array}{l}
0,1,2,0\\
1,1,0,1\\
0,1,0,1\\
1,0,1,1
\end{array}
\end{gathered}&
\begin{gathered}
p=17\\[3pt]
\begin{array}{l}
0,2,1,0\\
2,0,1,1\\
1,1,1,1\\
0,1,1,2
\end{array}
\end{gathered}&
\begin{gathered}
p=29\\[3pt]
\begin{array}{l}
2\mathpunct{.}\,3,0\mathpunct{.}\,2\\
1,1,2,3\\
2,1,2\mathpunct{.}\,1\\
1,2,3,1
\end{array}
\end{gathered}
\end{array}
\]

\SourceLeaf{89}%
\[
\begin{array}{c|c|c|c}
\begin{gathered}
p=37\\[3pt]
\begin{array}{l}
2,1,2,4\\
2,2,4,1\\
2,2,2,2\\
2,4,1,2
\end{array}
\end{gathered}&
\begin{gathered}
p=41\\[3pt]
\begin{array}{l}
0,4,3,2\\
4,2,2,2\\
3,2,3,2\\
2,2,2,4
\end{array}
\end{gathered}&
\begin{gathered}
p=53\\[3pt]
\begin{array}{l}
2,3,6,2\\
4,4,2,3\\
2,4,2,4\\
4,2,3,4
\end{array}
\end{gathered}&
\begin{gathered}
p=61\\[3pt]
\begin{array}{l}
4,3,2,6\\
3,3,6,3\\
4,3,4,3\\
3,6,3,3
\end{array}
\end{gathered}
\end{array}
\]

\[
\begin{array}{c|c|c}
\begin{gathered}
p=73\\[3pt]
\begin{array}{l}
5,6,4,2\\
6,2,5,5\\
4,5,4,5\\
2,5,5,6
\end{array}
\end{gathered}&
\begin{gathered}
p=89\\[3pt]
\begin{array}{l}
3,8,6,4\\
8,4,5,5\\
6,5,6,5\\
4,5,5,8
\end{array}
\end{gathered}&
\begin{gathered}
p=97\\[3pt]
\begin{array}{l}
2,6,7,8\\
6,8,5,5\\
7,5,7,5\\
8,5,5,6
\end{array}
\end{gathered}
\end{array}
\]

Da sich die Moduln der Formen $8n+1$ und $8n+5$ unterschiedlich verhalten, müssen wir sie getrennt behandeln. Wir beginnen mit den ersten.

\Needspace{14\baselineskip}\Article{16}
Das Zeichen $(00)$ gibt an, auf wie viele verschiedene Weisen die Gleichung $\alpha+1=\alpha'$ erfüllt werden kann, wobei $\alpha,\alpha'$ beliebige Zahlen aus $A$ bezeichnen. Für einen Modul der Form $8n+1$, wie wir ihn hier voraussetzen, gehören $\alpha'$ und $p-\alpha'$ zu derselben Sammlung. Daher können wir eleganter sagen, $(00)$ bezeichne die Anzahl der verschiedenen Möglichkeiten, die Gleichung $1+\alpha+\alpha'=p$ zu erfüllen. Offenbar kann an die Stelle dieser Gleichung auch die Kongruenz $1+\alpha+\alpha'\equiv0\ (\mathrm{mod.}\ p$ treten.

Ebenso
\begin{center}
\begin{minipage}{\linewidth}
bezeichnet $(01)$ die Anzahl der Lösungen der Kongruenz $1+\alpha+\beta\equiv0\ (\mathrm{mod.}\ p)$\par
$(02)$ die Anzahl der Lösungen der Kongruenz $1+\alpha+\gamma\equiv0$\par
$(03)$ die Anzahl der Lösungen der Kongruenz $1+\alpha+\delta\equiv0$\par
$(11)$ die Anzahl der Lösungen der Kongruenz $1+\beta+\beta'\equiv0$ usw.
\end{minipage}
\end{center}
Dabei bezeichnen $\beta$ und $\beta'$ beliebige Zahlen aus $B$, $\gamma$ Zahlen aus $C$ und $\delta$ Zahlen aus $D$. Daraus erhalten wir unmittelbar die folgenden sechs Gleichungen:
\[
\begin{gathered}
(01)=(10),\quad(02)=(20),\quad(03)=(30),\\
(12)=(21),\quad(13)=(31),\quad(23)=(32)
\end{gathered}
\]

Aus jeder gegebenen Lösung der Kongruenz $1+\alpha+\beta\equiv0$ erhält man eine Lösung von $1+\delta+\delta'\equiv0$, indem man für $\delta$ die Zahl zwischen $1\dotsfour p-1$\space%
\SourceLeaf{90}%
wählt, die $\beta\delta\equiv1$ erfüllt (sie gehört offenbar zu $D$), und für $\delta'$ den kleinsten positiven Rest des Produkts $\alpha\delta$ (der ebenfalls zu $D$ gehört). Ebenso ist der Rückweg von einer gegebenen Lösung von $1+\delta+\delta'\equiv0$ zu einer Lösung von $1+\alpha+\beta\equiv0$ möglich: Man wählt $\beta$ so, dass $\beta\delta\equiv1$ gilt, und setzt zugleich $\alpha\equiv\beta\delta'$. Beide Kongruenzen haben also gleich viele Lösungen, das heißt $(01)=(33)$.

Ebenso leiten wir aus der Kongruenz $1+\alpha+\gamma\equiv0$ die Kongruenz $\gamma'+\gamma''+1\equiv0$ her, wenn wir $\gamma'$ aus $C$ so wählen, dass $\gamma\gamma'\equiv1$ gilt, und $\gamma''$ aus derselben Sammlung kongruent zum Produkt $\alpha\gamma'$. Daraus folgt leicht, dass diese beiden Kongruenzen gleich viele Lösungen haben, also $(02)=(22)$.

Ebenso erhalten wir aus $1+\alpha+\delta\equiv0$ die Kongruenz $\beta+\beta'+1\equiv0$, wenn wir $\beta,\beta'$ so wählen, dass $\beta\delta\equiv1$, $\beta\alpha\equiv\beta'$ gilt. Damit ist $(03)=(11)$.

Schließlich gewinnen wir aus $1+\beta+\gamma\equiv0$ auf dieselbe Weise sowohl $\delta+1+\beta'\equiv0$ als auch $\gamma'+\delta'+1\equiv0$ und schließen daraus $(12)=(13)=(23)$.

Wir haben somit zwischen unseren sechzehn Unbekannten elf Gleichungen erhalten, sodass sich die Unbekannten auf fünf reduzieren und das Schema $S$ folgendermaßen geschrieben werden kann:
\[
\begin{array}{llll}
h,&i,&k,&l\\
i,&l,&m,&m\\
k,&m,&k,&m\\
l,&m,&m,&i
\end{array}
\]

Drei weitere Bedingungsgleichungen lassen sich leicht hinzufügen. Auf jede Zahl aus $A$, mit Ausnahme der letzten, $p-1$, muss eine Zahl aus einer der Sammlungen $A,B,C$ oder $D$ folgen. Also gilt
\[
(00)+(01)+(02)+(03)=2n-1
\]
und ebenso
\[
\begin{aligned}
(10)+(11)+(12)+(13)&=2n\\
(20)+(21)+(22)+(23)&=2n\\
(30)+(31)+(32)+(33)&=2n
\end{aligned}
\]
In den gerade eingeführten Zeichen ergeben die ersten drei Gleichungen:
\[
\begin{aligned}
h+i+k+l&=2n-1\\
i+l+2m&=2n\\
k+m&=n
\end{aligned}
\]

\SourceLeaf{91}%
Die vierte wird mit der zweiten identisch. Mithilfe dieser Gleichungen lassen sich drei Unbekannte eliminieren, womit alle sechzehn bereits auf zwei reduziert sind.

\Article{17}
Um sie vollständig zu bestimmen, untersuchen wir zweckmäßig die Anzahl der Lösungen der Kongruenz
\[
1+\alpha+\beta+\gamma\equiv0\ (\mathrm{mod.}\ p)
\]
wobei $\alpha,\beta,\gamma$ beliebige Zahlen aus $A,B,C$ bezeichnen. Offenbar ist der Wert $\alpha=p-1$ nicht zulässig, da $\beta+\gamma\equiv0$ nicht gelten kann. Setzt man für $\alpha$ nacheinander die übrigen Werte ein, so erhält man $h,i,k,l$ Werte von $1+\alpha$, die entsprechend zu $A,B,C,D$ gehören. Für jeden \emph{gegebenen} Wert von $1+\alpha$ in $A$, etwa $1+\alpha=\alpha^0$, hat die Kongruenz $\alpha^0+\beta+\gamma\equiv0$ so viele Lösungen wie $1+\beta'+\gamma'\equiv0$ (man setze $\beta\equiv\alpha^0\beta'$, $\gamma\equiv\alpha^0\gamma'$), also $(12)=m$ Lösungen. Ebenso hat für jeden gegebenen Wert von $1+\alpha$ in $B$, etwa $1+\alpha=\beta^0$, die Kongruenz $\beta^0+\beta+\gamma\equiv0$ so viele Lösungen wie $1+\alpha'+\beta'\equiv0$ (man setze $\beta\equiv\beta^0\alpha'$, $\gamma\equiv\beta^0\beta'$), also $(01)=i$ Lösungen. Entsprechend lässt sich für jeden gegebenen Wert von $1+\alpha$ in $C$, etwa $1+\alpha=\gamma^0$, die Kongruenz $\gamma^0+\beta+\gamma\equiv0$ auf so viele verschiedene Weisen lösen wie $1+\delta+\alpha'\equiv0$ (man setze $\beta\equiv\gamma^0\delta$, $\gamma\equiv\gamma^0\alpha'$); die Anzahl der Lösungen ist also $(03)=l$. Schließlich hat für jeden gegebenen Wert von $1+\alpha$ in $D$, etwa $1+\alpha=\delta^0$, die Kongruenz $\delta^0+\beta+\gamma\equiv0$ so viele Lösungen wie $1+\gamma'+\delta'\equiv0$ (man setze $\beta\equiv\delta^0\gamma'$, $\gamma\equiv\delta^0\delta'$), also $(23)=m$ Lösungen. Zählt man alles zusammen, so hat die Kongruenz $1+\alpha+\beta+\gamma\equiv0$ offenbar
\[
hm+ii+kl+lm
\]
verschiedene Lösungen.

\begingroup\rightskip=0pt plus1fil\relax Ganz ebenso erhalten wir, wenn für $\beta$ nacheinander die einzelnen Zahlen aus $B$ eingesetzt werden, dass die Summe $1+\beta$ entsprechend $(10),(11),(12),(13)$ oder $i,l,m,m$ Werte annimmt, die zu $A,B,C,D$ gehören. Für jeden \emph{gegebenen} Wert von $1+\beta$ aus diesen Sammlungen hat die Kongruenz $1+\beta+\alpha+\gamma\equiv0$ entsprechend $(02),(31),(20),(13)$ oder $k,m,k,m$ verschiedene Lösungen. Die Gesamtzahl aller Lösungen ist damit
\SourceLeaf{92}%
\[
=ik+lm+km+mm
\]
Zum selben Wert gelangen wir, wenn wir die Rechnung auf die Werte der Summe $1+\gamma$ gründen.\par\endgroup

\Article{18}
Aus diesen beiden Ausdrücken für dieselbe Anzahl ergibt sich die Gleichung
\[
0=hm+ii+kl-ik-km-mm
\]
und daraus durch Elimination von $h$ mithilfe von $h=2m-k-1$
\[
0=(k-m)^2+ii+kl-ik-kk-m
\]
Die letzten beiden Gleichungen aus Artikel 16 liefern $k=\frac12(l+i)$. Setzt man diesen Wert ein, so wird $ii+kl-ik-kk$ zu $\frac14(l-i)^2$, und die vorige Gleichung geht nach Multiplikation mit $4$ über in
\[
0=4(k-m)^2+(l-i)^2-4m
\]
Da $4m=2(k+m)-2(k-m)=2n-2(k-m)$ gilt, folgt
\[
2n=4(k-m)^2+2(k-m)+(l-i)^2
\]
oder
\[
8n+1=(4(k-m)+1)^2+4(l-i)^2
\]
Setzen wir also
\[
4(k-m)+1=a,\quad2l-2i=b
\]
so haben wir
\[
p=aa+bb
\]

Bekanntlich lässt sich $p$ auf genau eine Weise in zwei Quadrate zerlegen, wobei das ungerade als $aa$ und das gerade als $bb$ zu nehmen ist. Damit sind $aa,bb$ vollständig bestimmte Zahlen. Aber auch $a$ selbst ist vollständig bestimmt: Die Quadratwurzel ist positiv oder negativ zu nehmen, je nachdem die positive Wurzel die Form $4M+1$ oder $4M+3$ hat. Über die Bestimmung des Vorzeichens von $b$ werden wir gleich sprechen.

\Needspace{14\baselineskip}Verbindet man diese neuen Gleichungen mit den drei letzten Gleichungen aus Artikel 16, so sind die fünf Zahlen $h,i,k,l,m$ durch $a,b$ und $n$ folgendermaßen vollständig bestimmt:
\SourceLeaf{93}%
\[
\begin{aligned}
8h&=4n-3a-5\\
8i&=4n+a-2b-1\\
8k&=4n+a-1\\
8l&=4n+a+2b-1\\
8m&=4n-a+1
\end{aligned}
\]

Wollen wir anstelle von $n$ lieber den Modul $p$ einführen, so hat das Schema $S$, dessen einzelne Einträge zur Vermeidung von Brüchen mit $16$ multipliziert sind, folgende Gestalt:
\[
\begin{array}{l|l|l|l}
p-6a-11&p+2a-4b-3&p+2a-3&p+2a+4b-3\\
p+2a-4b-3&p+2a+4b-3&p-2a+1&p-2a+1\\
p+2a-3&p-2a+1&p+2a-3&p-2a+1\\
p+2a+4b-3&p-2a+1&p-2a+1&p+2a-4b-3
\end{array}
\]

\Article{19}\begingroup\postdisplaypenalty=10000\relax
Es bleibt zu zeigen, wie man das Vorzeichen für $b$ festlegt. In Artikel 10 haben wir bereits bemerkt, dass die Unterscheidung zwischen $B$ und $D$, die an sich nicht wesentlich ist, von der Wahl der Zahl $f$ abhängt. Für diese muss eine der beiden Wurzeln der Kongruenz $xx\equiv-1$ genommen werden, und die beiden Sammlungen werden vertauscht, wenn man statt der einen Wurzel die andere wählt. Das gerade angegebene Schema zeigt, dass eine entsprechende Vertauschung mit einem Wechsel des Vorzeichens von $b$ verbunden ist. Es lässt sich daher vorhersehen, dass zwischen dem Vorzeichen von $b$ und der Zahl $f$ ein Zusammenhang bestehen muss. Um ihn zu finden, bemerken wir zunächst: Bezeichnet $\mu$ eine nichtnegative ganze Zahl und nimmt man für $z$ alle Zahlen $1,2,3\dotsfour p-1$, so gilt modulo $p$ entweder $\Sigma z^\mu\equiv0$ oder $\Sigma z^\mu\equiv-1$, je nachdem $\mu$ nicht durch $p-1$ teilbar oder dadurch teilbar ist. Der letzte Teil des Satzes folgt daraus, dass für einen Wert von $\mu$, der durch $p-1$ teilbar ist, die Beziehung $z^\mu\equiv1$ gilt; den ersten Teil beweisen wir folgendermaßen. Bezeichnet $g$ eine primitive Wurzel, so stimmen alle $z$ mit den kleinsten Resten aller $g^y$ überein, wenn man für $y$ alle Zahlen $0,1,2,3\dotsfour p-2$ nimmt. Daher gilt $\Sigma z^\mu\equiv\Sigma g^{\mu\pdot y}$. Nun ist
\[
\begin{gathered}
\Sigma g^{\mu y}=\frac{g^{\mu(p-1)}-1}{g^\mu-1},\quad\text{also}\\
(g^\mu-1)\Sigma z^\mu\equiv g^{\mu(p-1)}-1\equiv0
\end{gathered}
\]
Für einen Wert von $\mu$, der nicht durch $p-1$ teilbar ist, kann $g^\mu$ nicht $1$ kongruent sein; das heißt, $g^\mu-1$ kann nicht durch $p$ teilbar sein. Daraus folgt $\Sigma z^\mu\equiv0$. Q. E. D.\par\endgroup

\SourceLeaf{94}%
Entwickelt man nun die Potenz $(z^4+1)^{\frac14(p-1)}$ nach dem Binomialtheorem, so ergibt der vorige Hilfssatz
\[
\Sigma(z^4+1)^{\frac14(p-1)}\equiv-2\ (\mathrm{mod.}\ p)
\]
Die kleinsten Reste aller $z^4$ liefern aber sämtliche Zahlen in $A$, wobei jede viermal vorkommt. Unter den kleinsten Resten von $z^4+1$ haben wir daher
\[
\begin{array}{rl}
4(00)&\text{zu }A\\
4(01)&\text{zu }B\\
4(02)&\text{zu }C\\
4(03)&\text{zu }D
\end{array}
\]
zu diesen Klassen gehörige, und vier sind $=0$ (nämlich für $z^4\equiv p-1$). Aus den Kriterien für die Sammlungen $A,B,C,D$ folgt
\[
\begin{aligned}
\Sigma(z^4+1)^{\frac14(p-1)}\equiv{}&4(00)+4f(01)\\
&-4(02)-4f(03)
\end{aligned}
\]
und damit
\[
-2\equiv4(00)+4f(01)-4(02)-4f(03)
\]
oder, wenn man für $(00),(01)$ usw. die im vorigen Artikel gefundenen Werte einsetzt,
\[
-2\equiv-2a-2-2bf
\]
Wir schließen also, dass stets $a+bf\equiv0$ gelten muss, oder nach Multiplikation mit $f$
\[
b\equiv af
\]
Diese Kongruenz bestimmt das Vorzeichen von $b$, wenn die Zahl $f$ bereits gewählt ist, oder die Zahl $f$, wenn das Vorzeichen von $b$ anderweitig vorgeschrieben wird.
